\documentclass[../main.tex]{subfiles}

\begin{document}

\section{Implementation: Key Code Excerpts}\label{sec:appendix_code}

The full MATLAB sources (\texttt{main\_task1.m}--\texttt{main\_task4.m} and \texttt{ode\_burn.m}) live in the \texttt{HM1/} folder of the repository. Here I collect only the three routines that do the actual work: the powered-flight right-hand side, the shooting residual, and the continuation sweep. The four tasks reuse the same building blocks, so showing them once is enough.

The optimal dynamics are integrated by the right-hand side in Listing~\ref{lst:oderhs}. Four of the five costates are written in closed form (two constants and the linear ramp $\lambda_{v_y}$), so only $\lambda_m$ is actually propagated, and the thrust angle is read off the velocity costate with a single \texttt{atan2}.

\lstinputlisting[style=matlab,
  caption={Powered-flight right-hand side with linear-tangent steering (\texttt{ode\_burn.m}).},
  label=lst:oderhs]{../ode_burn.m}

Listing~\ref{lst:shooting} is the four-residual shooting function shared by Tasks~1, 2 and~4. The two numerical refinements discussed in Section~\ref{sec:shooting} are both visible: the state is initialized with $\lambda_{m,0}=1$ (the costate normalization), and the free-time condition $\mathcal{H}(0)=0$ is imposed algebraically at $t_0$ rather than at $t_f$.

\lstinputlisting[style=matlab, firstnumber=227, firstline=227, lastline=273,
  caption={Shooting residual for a single burn arc (\texttt{main\_task1.m}).},
  label=lst:shooting]{../main_task1.m}

The parameter sweeps depend on continuation: each new value of $Q$ is solved starting from the previous converged costate vector, sweeping both forward and backward from a seed point near $Q=3$. Listing~\ref{lst:continuation} shows the two sweeps.

\lstinputlisting[style=matlab, firstnumber=65, firstline=65, lastline=93,
  caption={Continuation sweep in $Q$ (excerpt from the per-altitude loop, \texttt{main\_task1.m}).},
  label=lst:continuation]{../main_task1.m}

\end{document}
